\chapter{Introduction}

Object detection is a core problem in computer vision with applications in autonomous systems, intelligent surveillance, robotics, and medical image analysis. In recent years, deep convolutional detectors have achieved strong accuracy on large-scale benchmarks such as MS COCO, driven by advances in backbone design, multi-scale feature fusion, and improved training objectives \cite{lin2014coco,he2016resnet,lin2017fpn,lin2017focal}. However, the detectors that obtain the highest accuracy often rely on deep backbones and heavy feature extractors, which makes them expensive to deploy in settings with limited memory, latency, or energy budgets. This creates a practical tension: modern detection systems must remain accurate while also becoming efficient enough for real-world inference constraints. More broadly, efficient deployment has also motivated classical model-compression methods such as pruning, quantization, and coding-based compression pipelines \cite{han2015learning,hubara2016quantized,han2016deep}.

One-stage detectors are especially attractive in this context because they offer a clean trade-off between accuracy and speed. RetinaNet is a particularly influential one-stage detector because it combines a Feature Pyramid Network (FPN) with focal loss to address extreme foreground-background imbalance during dense prediction \cite{lin2017fpn,lin2017focal}. Even so, RetinaNet performance depends strongly on the backbone capacity. A detector with a deep \texttt{ResNet-101} feature extractor typically yields better representations than a detector with a lighter \texttt{ResNet-18} backbone, but the latter is preferable when deployment efficiency matters. A natural question is therefore whether the representational strength of a larger RetinaNet can be transferred into a smaller one without reproducing the full computational cost at inference time.

\ArchitectureFigure{Overview of the teacher-student RetinaNet distillation architecture used throughout this dissertation.}{fig:architecture-overview}

\section{Background and Context}

Knowledge distillation offers a principled framework for this goal. In the standard teacher-student setting, a large model is first used as a source of supervision and a smaller student is trained to imitate selected aspects of the teacher while still optimizing the task objective \cite{hinton2015distilling}. Earlier work explored distillation through softened logits, intermediate hints, and attention transfer \cite{romero2015fitnets,zagoruyko2017paying}. In object detection, distillation is more challenging because the model must preserve classification, localization, and multi-scale feature reasoning simultaneously. Recent detection-oriented distillation methods therefore supervise feature maps, spatial relations, and foreground regions rather than relying only on final predictions \cite{wang2019distilling,dai2021gid,guo2021defeat,kang2021icd,du2021frs,li2022rank,yang2022focal}. The continued importance of distillation has even motivated recent work on model provenance, where Shi et al.\ study how to detect whether an open-weights student has been distilled from a teacher under restricted access assumptions \cite{shi2025kddetection}.

This dissertation studies knowledge distillation for efficient one-stage object detection using RetinaNet as the common detection framework. The implemented system uses a frozen teacher detector with a \texttt{ResNet-101} or \texttt{ResNet-50} FPN backbone and trains a smaller student detector, including a \texttt{ResNet-18} configuration, on the same detection task. In addition to the native RetinaNet detection losses, the student is optimized with a feature-based distillation loss applied to FPN representations. Inspired by recent structural knowledge distillation for object detection \cite{derijk2022structural}, the method replaces pointwise feature matching with a structural similarity objective that encourages the student to reproduce the teacher's local spatial organization within the feature pyramid.

The central idea is that feature distillation for detection should preserve not only activation magnitude but also structural relationships that matter for object localization across scales. Structural similarity, originally introduced for image quality assessment, is well suited to this purpose because it compares local statistics such as mean, variance, and correlation rather than treating each activation independently \cite{wang2004ssim}. Building on the SSIM-based distillation perspective of de Rijk et al.\ \cite{derijk2022structural}, this dissertation emphasizes explicit hierarchy-aware supervision over RetinaNet FPN levels so that the student learns teacher structure at multiple resolutions while still being trained end-to-end for dense detection.

\section{Objectives}

The objectives of this dissertation are as follows:

\begin{itemize}
    \item To develop a practical teacher-student distillation framework for RetinaNet with interchangeable ResNet-FPN teacher and student backbones.
    \item To adapt SSIM-based structural distillation to level-aligned FPN features so that multi-scale spatial structure can be transferred directly between detectors.
    \item To analyse the dimensional compatibility of teacher and student features and justify a post-FPN distillation interface that does not require a separate learned adaptation layer.
    \item To implement an evaluation protocol for comparing student-only RetinaNet training with teacher-guided training under COCO-style metrics.
\end{itemize}

\section{Achievements}
% TODO 
The work completed so far has established the main conceptual and technical foundations of the project. A structured review of the literature has been written to position the dissertation within modern detector distillation research. A detailed teacher-student architecture has been formulated for RetinaNet, together with a combined loss that couples standard detection supervision to an SSIM-based feature distillation term. The dissertation also now documents the dimensionality of the teacher and student feature spaces and explains why distillation after the FPN allows direct comparison without an additional adaptation layer. Finally, the training and evaluation protocol has been specified clearly enough to support the next experimental stage of the work.

\section{Practical Considerations and Constraints}

The scope of this dissertation is shaped by a practical but important hardware constraint: experimental training is limited to access to a single rented NVIDIA Tesla V100 GPU with 32 GB of memory, rather than to a dedicated multi-GPU research server \cite{nvidia2017v100}. Although the V100 remains a capable accelerator for deep learning workloads, working within a one-GPU rental budget affects every stage of the study, including feasible image resolution, batch size, training duration, and the number of controlled ablation runs that can be completed within the dissertation timeline. These limits are especially relevant for object detection, where dense prediction, feature pyramids, large backbone activations, and multi-scale supervision collectively create a much heavier memory and compute footprint than standard image classification settings.

Within this context, the choice to study teacher-student distillation with RetinaNet is deliberate rather than merely conservative. RetinaNet provides a strong and well-established one-stage baseline, while its ResNet-FPN design offers a comparatively stable experimental platform for analyzing distillation behaviour \cite{lin2017fpn,lin2017focal}. By keeping the detector family fixed and varying the backbone capacity between teacher and student, the dissertation can focus on the main research question: whether structural feature supervision can transfer useful knowledge from a larger detector to a smaller one under realistic deployment constraints. A more heterogeneous study mixing multiple detector paradigms would make it harder to attribute any observed performance differences specifically to the distillation objective.

This compute budget also limits the ability to explore stronger attention-based detectors in a rigorous way. Transformer-based detection has become an important direction in the literature, with DETR showing that object detection can be formulated through end-to-end set prediction with attention mechanisms \cite{carion2020detr}, and hierarchical transformer backbones such as Swin Transformer becoming competitive general-purpose representations for dense vision tasks \cite{liu2021swin}. These architectures are highly relevant to the broader field, but incorporating them properly would require more than simply swapping one backbone for another. A fair comparison would involve retraining or carefully fine-tuning additional teacher and student configurations, tuning distinct optimization schedules, and repeating the same ablation process across a substantially larger design space. Under a single rented V100, that expansion would consume a large fraction of the available project time and budget before the proposed distillation mechanism itself could be evaluated with sufficient care.

The same constraint applies even more strongly to recent foundation-style detection models. Systems such as GLIP and Grounding DINO combine object detection with language grounding and large-scale multimodal pre-training, thereby moving detection toward open-vocabulary and open-set reasoning \cite{li2022glip,liu2023groundingdino}. GLIP, for example, is pre-trained on 27 million grounding examples, including millions of web-crawled image-text pairs \cite{li2022glip}, while Grounding DINO extends this trend by coupling transformer-based detection with grounded pre-training for arbitrary text-specified categories \cite{liu2023groundingdino}. These models are scientifically important because they greatly broaden what a detector can represent, but reproducing their training pipelines, or even conducting fair fine-tuning studies, would move the dissertation into a very different compute regime. They also introduce extra experimental variables such as language encoders, prompt construction, and pre-training data provenance, which would dilute the present focus on controlled feature distillation for closed-set object detection.

Instead tring to greatly advance SOTA with a more complex architecture, this dissertation focuses on efficiency of smaller detectors through distillation, which is a more tractable and still highly relevant research question under the given constraints. The proposed method can be implemented and evaluated within the available compute budget, while still contributing to the important problem of how to make modern detection systems more efficient for real-world deployment.


% TODO

\section{Overview of Dissertation}

Chapter 2 reviews the background theory and related literature in deep object detection, knowledge distillation, detector-specific distillation, and structural similarity. Chapter 3 presents the proposed methodology, including the problem formulation, the teacher-student architecture, and the dimensional analysis of the feature interface. Chapter 4 details the training objective, FPN feature extraction process, SSIM-based distillation mechanism, and evaluation protocol. Chapter 5 concludes the current draft by evaluating the progress made so far and outlining the most important next steps for the experimental phase.
